%% Generated by verification/make_figures.py from regular-graph-harmonic-quotient-examples.tex.
%% Figure fig:harmonic-quotient_d4reg12a+v9e12ds444222222 (replaces figures/harmonic-quotient_d4reg12a+v9e12ds444222222-G1.tex).
%% Gamma_1 = image of phi: 9 vertices, 12 edges.
%% Vertex y is filled with color hq<y>.
%% Layout: force-directed (1 crossings).
\begin{tikzpicture}[hqgraph]
\node[hqv, fill=hq0] (v0) at (1.770,0.146) {0};
\node[hqv, fill=hq1] (v1) at (0.747,0.795) {1};
\node[hqv, fill=hq2] (v2) at (-1.770,0.149) {2};
\node[hqv, fill=hq3] (v3) at (1.646,-0.771) {3};
\node[hqv, fill=hq4] (v4) at (-0.463,0.258) {4};
\node[hqv, fill=hq5] (v5) at (-1.335,-1.587) {5};
\node[hqv, fill=hq6] (v6) at (-1.147,-0.543) {6};
\node[hqv, fill=hq7] (v7) at (0.358,-1.120) {7};
\node[hqv, fill=hq8] (v8) at (-0.888,1.587) {8};
\draw[hqe] (v0) -- (v1);
\draw[hqe] (v0) -- (v7);
\draw[hqe] (v1) -- (v3);
\draw[hqe] (v1) -- (v4);
\draw[hqe] (v1) -- (v8);
\draw[hqe] (v2) -- (v4);
\draw[hqe] (v2) -- (v5);
\draw[hqe] (v2) -- (v6);
\draw[hqe] (v2) -- (v8);
\draw[hqe] (v3) -- (v7);
\draw[hqe] (v5) -- (v7);
\draw[hqe] (v6) -- (v7);
\end{tikzpicture}%
